% ==============================================================================
% ssec:cs_srovnani --- Průřezové srovnání případů
% Zdroje: PennyMarket_KS2026, KS_CeskaPosta_2025, KSVS_OSPZV_ASO_2026,
%         KSVS_OSPPP_2026, CMKOS_ZpravaKV2025, CMKOS_PrinosKV2025, zakon_zp_2006
% ==============================================================================

lze nejsnáze provést ve třech dimenzích -- mzdové základny, výše nárokových příplatků a~rozsahu nadstandardních benefitů.
Tab.~\ref{tab:cs_srovnani} sestavuje tyto parametry vedle sebe. Všechny hodnoty jsou převzaty z~textů smluv s~účinností pro rok 2026 a~doplněny zákonným minimem podle \aclgen{ZP} a~prováděcích nařízení.
Vedle parametrů samotných je však rozhodující i~to, co za~ně každá strana směnila, neboť právě bilance motivací vysvětluje rozdíly mezi sektory.\par

\begin{table}[H]
  \centering
  \caption[Průřezové srovnání analyzovaných kolektivních smluv]{Průřezové srovnání čtyř analyzovaných kolektivních smluv, \aca{geo-CZ}, 2026. Zdroje dat: \acs{PKS}~\cite{PennyMarket_KS2026}~\cite{KS_CeskaPosta_2025}, \acs{KSVS}~\cite{KSVS_OSPZV_ASO_2026}~\cite{KSVS_OSPPP_2026} a~\acs{ZP}~\cite{zakon_zp_2006}}
  \label{tab:cs_srovnani}
  \begin{tabularx}{\linewidth}{Xrrrrr}
    \toprule
    \textbf{Parametr} & \textbf{\acs{ZP} min.} & \textbf{Penny} & \textbf{\acs{ČP}} & \textbf{zeměděl.} & \textbf{finance} \\
    \midrule
    Min.~mzda [Kč/měs.] & 22\,400 & 22\,400 & 22\,400 & 22\,400 & 28\,000 \\
    Přípl.~přesčas [\% \acs{PV}] & 25 & 25 & 25 & 25 & 25 \\
    Přípl.~noc [\% \acs{PV}] & 10 & 10 & 15 & 20 & 10 \\
    Přípl.~so/ne [\% \acs{PV}] & 10 & 10 & 12 & 20 & 50 \\
    Vyrovnávací období [týd.] & 26 & 52 & 26 & 52 & 26 \\
    Dovolená [týd.] & 4 & 6 & 5 & 4+ & 5 \\
    Sick days [dny] & 0 & 0 & 2 & 0 & 3 \\
    Příspěvek \acs{DPS} [Kč/měs.] & 0 & 800 & 500 & až 1\,350 & --- \\
    Sociální fond [\% mezd] & --- & --- & 1 & až~1 & 2{,}5 \\
    Účinnost [roky] & --- & 1 & 1+prol. & 1 & 2 \\
    Počet \acs{OO} & --- & 5 & 9 & 1 & 1 \\
    \bottomrule
  \end{tabularx}
\end{table}

Z~tabulky vyplývají tři pozorování s~přesahem k~obecné hypotéze této práce:
\begin{enumerate}
  \item samotná existence kolektivní smlouvy negarantuje silnou ochranu práce. \Acf{PKS} Penny~Market často kopíruje zákonná minima a~výměnou za~flexibilitu vyrovnávacího období (\SI{52}{\week}) nabízí pouze fiskálně nenáročné benefity. Naopak \ac{ČP} a~obě \acp{KSVS} překračují zákonné příplatky až pětinásobně;
  \item smluvní nadstandardy kopírují sektorové motivace a~tržní postavení aktérů. V~sektorech s~vysokou přidanou hodnotou (finance) strany konstruují komplexní ochranu s~vysokými stabilizačními příplatky. Odvětvová smlouva v~zemědělství pak slouží organizovaným zaměstnavatelům k~eliminaci mzdového dumpingu;
  \item komparace prokazuje nezastupitelnost odvětvového vyjednávání. Podniková ujednání nedokáží kompenzovat chybějící sektorový standard v~odvětvích s~vysokou fragmentací a~nízkou odborovou organizovaností, jako je maloobchod.
\end{enumerate}\par

Typické znaky vyjednávání v~nízkopříjmových sektorech je detailně sjednávaná \acs{BOZP} (kontroly, psychosociální prevence) a~pasivní přístup ke~vzdělávání, což odpovídá tradičnímu defenzivnímu zaměření českého vyjednávání.~\cite{Eurofound_CBBeyondPay_2024}\par

Komparace odvětvových smluv odhaluje limity českého sociálního dialogu. Zatímco \ac{KSVS} v~zemědělství plošně chrání sektor díky zákonné extenzi, v~bankovnictví k~extenzi nedošlo. Smlouva je~tak závazná pouze pro~členy \acgen{SBP}, kteří však reprezentují \num{16}~klíčových zaměstnavatelů s~více než \SI{46000}{\person}. To tvoří významnou část trhu práce, a~tak efektivní standard v~odvětví. V~maloobchodě (\ac{NACE}~47.11) naopak fragmentace zaměstnavatelů a~marginální odborová organizovanost vznik jakýchkoli odvětvových minim v~duchu evropské směrnice znemožňuje.~\cite{eu_directive_2022_2041}~\cite{CMKOS_ZpravaKV2025}\par
